% =====================================================================
\chapter{Experimental Setup}
\label{ch:experimental-setup}
% =====================================================================

This chapter specifies the datasets (\S\ref{sec:exp-datasets}), the continual
task construction (\S\ref{sec:exp-tasks}), the model and training configuration
(\S\ref{sec:exp-config}), the baselines (\S\ref{sec:exp-baselines}), and the
evaluation methodology (\S\ref{sec:exp-eval}). Together these define the FCS-CL
benchmark used in Chapter~\ref{ch:results}.

\section{Datasets}
\label{sec:exp-datasets}

We use three public visually-rich document datasets that together span forms and
receipts and a wide range of label-space sizes. Their statistics are summarised in
Table~\ref{tab:dataset-stats}.

\begin{description}
  \item[FUNSD] noisy scanned forms with three entity types (HEADER, QUESTION,
        ANSWER) plus background, giving 7 BIO tags. Small and noisy, it is the
        clearest test of the KEY/VALUE relation.
  \item[CORD] receipt understanding with 30 fine-grained classes grouped into five
        super-classes (menu, sub\_total, total, void\_menu, sub), giving 61 BIO
        tags. Its fine-grained label space makes it the natural host for the
        class-incremental scenario.
  \item[SROIE] scanned receipts with four fields (company, address, date, total),
        giving 9 BIO tags.
\end{description}

\begin{table}[t]
  \centering
  \caption[Dataset statistics]{Statistics of the three datasets. Tag counts are
  BIO-aware and include the background ``O'' tag. Split sizes follow the standard
  releases used in this work.}
  \label{tab:dataset-stats}
  \begin{tabular}{lcccl}
    \toprule
    \textbf{Dataset} & \textbf{Train} & \textbf{Test} & \textbf{BIO tags} & \textbf{Domain} \\
    \midrule
    FUNSD & 149 & 50  & 7  & scanned forms \\
    CORD  & 800 & 100 & 61 & receipts (fine-grained) \\
    SROIE & 626 & 347 & 9  & scanned receipts \\
    \bottomrule
  \end{tabular}
\end{table}

\section{Continual-Learning Task Construction}
\label{sec:exp-tasks}

Three scenarios are constructed over these datasets.

\paragraph{Class-incremental (CIL-CORD).} CORD's 30 fine-grained classes are split
into five sessions of six classes each. The partition follows super-class
boundaries where possible, ordered as a curriculum from core menu fields to
auxiliary void/sub categories, with sessions balanced to roughly equal training
size. Each session tags its super-class context as metadata so that forgetting can
be analysed within versus across super-classes.

\paragraph{Domain-incremental (DIL).} FUNSD, SROIE, and CORD are remapped onto a
unified four-entity schema (HEADER, KEY, VALUE, OTHER; 9 BIO tags) and presented
in order of increasing distance from that canonical schema:
FUNSD\,$\rightarrow$\,SROIE\,$\rightarrow$\,CORD. FUNSD's QUESTION/ANSWER maps to
KEY/VALUE; SROIE's COMPANY maps to KEY and its other fields to VALUE; CORD's
productive super-classes map to VALUE and its void/sub categories to OTHER. The
mapping and its trade-offs (notably the sparsity of KEY outside FUNSD) are recorded
as part of the benchmark.

\paragraph{Mixed.} A six-session scenario interleaves class-incremental and domain
shifts to stress methods under combined non-stationarity.

For robustness, the diagnosis is additionally repeated under at least one alternate
task order, and all scenarios are run for three seeds.

\section{Model Architecture and Training Configuration}
\label{sec:exp-config}

The backbone is LayoutLMv3-base (\S\ref{sec:impl-framework}) with a token-classification
head expanded as new labels appear. Unless otherwise stated, training uses the
pilot defaults---learning rate $5\times10^{-5}$, batch size 8, and a fixed number
of epochs per task---with AdamW. Hyperparameter sensitivity is treated as a
separate ablation in the main grid rather than tuned per method, to keep the
comparison fair. Determinism is enforced via fixed seeds and deterministic cuDNN.
Full hyperparameters are listed in Appendix~\ref{app:hyperparameters}.

\section{Baselines and Continual-Learning Strategies}
\label{sec:exp-baselines}

We compare against ten baselines spanning the major families, plus the selected
DocCL method, all under the shared interface of \S\ref{sec:impl-cl-core}:

\begin{itemize}
  \item \textbf{Bounds:} \emph{Naive} sequential fine-tuning (lower bound) and
        \emph{Joint} multi-task training on all data (oracle upper bound).
  \item \textbf{Regularisation:} EWC~\citep{kirkpatrick2017ewc},
        LwF~\citep{li2017lwf}.
  \item \textbf{Replay:} ER~\citep{rolnick2019er}, DER++~\citep{buzzega2020der}.
  \item \textbf{Prompt-based:} L2P~\citep{wang2022l2p},
        DualPrompt~\citep{wang2022dualprompt}, CODA-Prompt~\citep{smith2023coda}.
  \item \textbf{LoRA-based:} O-LoRA~\citep{wang2023olora}.
\end{itemize}

Methods are reproduced and validated against their published behaviour before being
extended (for example, ER is made to work before DER++). The selected DocCL method
is added once the diagnosis fixes its design (\S\ref{sec:method-selected}).

\section{Evaluation Methodology}
\label{sec:exp-eval}

\subsection{Validation Protocol}
\label{sec:eval-protocol}

After each task the model is evaluated on all tasks seen so far, populating the
accuracy matrix $R$. The primary task metric is entity-level F1 with BIO-aware
span matching (\texttt{seqeval}).

\subsection{Metrics}
\label{sec:eval-metrics}

From the accuracy matrix $R[i,j]$ (F1 on task $j$ after training task $i$) we
report the standard continual-learning metrics~\citep{lopezpaz2017gem}:

\begin{align}
  \mathrm{AA}  &= \frac{1}{T}\sum_{j=1}^{T} R[T,j] && \text{(average accuracy)},\\
  \mathrm{BWT} &= \frac{1}{T-1}\sum_{j=1}^{T-1}\big(R[T,j]-R[j,j]\big) && \text{(backward transfer)},\\
  \mathrm{FWT} &= \frac{1}{T-1}\sum_{j=2}^{T}\big(R[j-1,j]-\bar{b}_j\big) && \text{(forward transfer)},
\end{align}
where $\bar{b}_j$ is a random-initialisation reference for task $j$, and average
forgetting is $\mathrm{AF}=-\mathrm{BWT}$. In addition we report the
\emph{per-component} drift metrics of \S\ref{sec:method-diag-metrics} (per-layer
CKA and per-group Fisher), which are the diagnostic contribution of this thesis.

\subsection{Repeated-Seed Protocol and Statistical Testing}
\label{sec:eval-seeds}

All main results are reported as mean\,$\pm$\,standard deviation over three seeds
(42, 123, 7). Comparisons between the selected method and the strongest baseline
use a paired test (paired $t$-test or Wilcoxon signed-rank), and the per-component
diagnosis uses a Mann--Whitney U test with Bonferroni correction across the eight
parameter groups (\S\ref{sec:method-diag-hypotheses}).

\section{Summary}
\label{sec:exp-summary}

The FCS-CL benchmark comprises three datasets, three continual scenarios, ten
baselines plus the selected method, and a dual evaluation: standard
continual-learning metrics (AA, BWT, FWT, AF) together with the per-component
diagnostic metrics. All results are reported over three seeds with statistical
tests. The next chapter presents the findings.
